\exercisesection{Stellar Algebras}

\begin{enumerate}
\def\labelenumi{\arabic{enumi})}
\item
  Find a unital involutive Banach algebra A and a unitary element \(u \in A\) such that \(\|u\| \neq 1\) .
\item
  Let A be a Banach algebra and \(A_{1} \subset A\) a closed subalgebra. If \(A_{1}\) is an involutive algebra, then every character of \(A_{1}\) extends to a character of A.
\item
  Let A be a unital involutive algebra. Suppose that for every x in A, the spectrum of x does not contain -1.
\end{enumerate}

\begin{enumerate}
\def\labelenumi{\alph{enumi})}
\item
  For every idempotent element x of A, show that there exists a hermitian idempotent element y of A such that \(xA = yA\) . (Consider the inverse z of \(1 + (x - x^{*})(x^{*} - x)\) ; show that z commutes with x and set \(y = xx^{*}z\) .)
\item
  For every idempotent element \(x\) of A, show that there exists a hermitian idempotent element \(y\) of A and \(u\) invertible in A such that \(x = uyu^{-1}\) . (If \(x\mathrm{A} = y\mathrm{A}\) , set \(u = (1 - (y - x))^{-1}\) .)
\end{enumerate}

\begin{enumerate}
\def\labelenumi{\arabic{enumi})}
\setcounter{enumi}{3}
\item
  Let A be a normed algebra, equipped with a continuous involution. If one sets \(\|x\|' = \sup(\|x\|, \|x^{*}\|)\) for all \(x \in A\) , A becomes an involutive normed algebra, and the new norm is equivalent to the old one.
\item
  Let A be a commutative Banach algebra without radical. Every involution of A is continuous. (Use prop. 9 of I, p.~40.)
\end{enumerate}

\begin{enumerate}
\def\labelenumi{¶ \arabic{enumi})}
\setcounter{enumi}{5}
\tightlist
\item
  Let A be a commutative unital involutive Banach algebra whose every character is hermitian. Suppose that, for every \(x \in A\) non-invertible hermitian element, one has \(\|(x - \lambda)^{-1}\| = o(\mathcal{I}(\lambda)^{-2})\) when \(\lambda\) tends to 0 with \(\mathcal{I}(\lambda) \neq 0\) . Then every closed ideal I of A is an intersection of maximal ideals. (Let \(\varphi: A \longrightarrow A/I\) be the canonical morphism. Let \(x \in A\) such that \(\varphi(x)\) be in the radical of A/I. It must be proved that \(\varphi(x) = 0\) . Reduce to the case where \(x\) is hermitian. Then use exerc. 29 of I, p.~161).
\end{enumerate}

\begin{enumerate}
\def\labelenumi{\arabic{enumi})}
\setcounter{enumi}{6}
\item
  Let A be a commutative unital involutive Banach algebra. For every character of A to be hermitian, it is necessary and sufficient that \(1 + xx^{*}\) be invertible for every \(x \in A\) . (If every character of A is hermitian, \(\mathcal{G}(1 + xx^{*})\) is everywhere \textgreater{} 0 on \(\mathsf{X}(\mathsf{A})\) . If a character \(\chi\) of A is non-hermitian, there exists a \(x \in A\) hermitian element such that \(\chi(x) = i\) , whence \(\chi(1 + xx^{*}) = 0\) and \(1 + xx^{*}\) is not invertible.)
\item
  Construct an example of a group G such that the involutive Banach algebra \(\mathcal{M}^{1}(G)\) (I, p.~100, example 4) is not a star algebra.
\item
  Let A be a unital Banach algebra. The group of unitary elements of A is a maximal bounded subgroup of A*.
\item
  Let A be an involutive Banach algebra and \(\widetilde{\mathsf{A}}\) the algebra obtained by adjoining a unit element. Show that the norm defined in number I, p.~15 on \(\widetilde{\mathsf{A}}\) does not coincide, in general, with the norm defined by prop. 3 of I, p.~105.
\item
  Let A be an involutive Banach algebra, whose involution satisfies \(\|x^{*}x\| = \|x\| \cdot \|x^{*}\|\) for all \(x \in A\) . Then A is a star algebra. (One has \(\|y^{2}\| = \|y\|^{2}\) for hermitian y, whence \(\|y\| = \varrho(y)\) ; observing that \(\mathrm{Sp}_{\mathrm{A}}^{\prime}(x^{*}) = \overline{\mathrm{Sp}_{\mathrm{A}}^{\prime}(x)}\) for all \(x \in A\) , we therefore have
\end{enumerate}

\[
\| x ^ {*} \| \cdot \| x \| = \| x ^ {*} x \| = \varrho (x ^ {*} x) \leqslant \varrho (x ^ {*}) \varrho (x) = \varrho (x ^ {2}) \leqslant \| x \| ^ {2},
\]

whence \(\|x^{*}\|\leqslant\|x\|\) and \(\|x^{*}\|=\|x\|\) .

\begin{enumerate}
\def\labelenumi{\arabic{enumi})}
\setcounter{enumi}{11}
\item
  Let A be a star algebra, N the set of normal elements of A, \(x \in N\) and V a neighborhood of 0 in C. There exists a neighborhood U of x in N such that, for every \(y \in U\) , one has \(\mathrm{Sp}_{\mathrm{A}}(y) \subset \mathrm{Sp}_{\mathrm{A}}(x) + \mathrm{V}\) and \(\mathrm{Sp}_{\mathrm{A}}(x) \subset \mathrm{Sp}_{\mathrm{A}}(y) + \mathrm{V}\) . (Use the cor. of 1 of I, p.~108, prop. 3 of I, p.~23 and prop. 10 of I, p.~76.)
\item
  Let A be a unital Banach algebra over C, \(x \in A\) , and \(S = \mathrm{Sp}_{\mathrm{A}}(x)\) . Suppose that there exists a morphism \(\varphi\) of \(\mathcal{C}(\mathrm{S})\) into A which sends 1 to 1 and z to x, denoting by z the identity map of S. Then \(\|\varphi(f)\| \geqslant \|f\|\) for every \(f \in \mathcal{C}(\mathrm{S})\) . In particular, if the morphism \(\varphi\) is continuous, then \(\varphi\) is a homeomorphism onto its image. (One may assume A commutative. Let \(\lambda \in S\) . There exists \(\chi \in \mathrm{X}(\mathrm{A})\) such that \(\lambda = \chi(x) = \chi(\varphi(z)) = (\mathrm{X}(\varphi)(\chi))(z)\) . Therefore \(\mathrm{X}(\varphi)(\chi)\) be the character \(f \mapsto f(\lambda)\) of \(\mathcal{C}(\mathrm{S})\) . Let \(f \in \mathcal{C}(\mathrm{S})\) . According to what precedes, there exists \(\chi \in \mathrm{X}(\mathrm{A})\) such that \(|f|\) attains its maximum at \(\mathrm{X}(\varphi)(\chi)\) . Then \(|\chi(\varphi(f))| = \|f\|\) , hence \(\|\varphi(f)\| \geqslant \|f\|\) .)
\item
  Let G be a locally compact group. If Stell(G) has an identity element, then the group G is discrete. (Use the fact that, in the Hilbert space \(L^{2}(G)\) , the identity map is the norm limit of endomorphisms \(\gamma_{f}\) , where \(\gamma\) is the left regular representation and where \(f \in L^{1}(G)\) .)
\item
  \begin{enumerate}
  \def\labelenumii{\alph{enumii})}
  \tightlist
  \item
    Let A be a commutative stellar algebra. If x and y are positive elements of A such that \(x \leqslant y\) , then \(x^{2} \leqslant y^{2}\) .
  \end{enumerate}
\end{enumerate}

\begin{enumerate}
\def\labelenumi{\alph{enumi})}
\setcounter{enumi}{1}
\item
  There exists a stellar algebra A and elements x and y of A such that \(0 \leqslant x \leqslant y\) but such that \(x^{2}\) is not \(\leqslant y^{2}\) .
\item
  There exists a stellar algebra A and elements \(x\) and \(y\) in A such that \(|x + y|\) is not \(\leqslant |x| + |y|\) .
\end{enumerate}

\begin{enumerate}
\def\labelenumi{\arabic{enumi})}
\setcounter{enumi}{15}
\tightlist
\item
  Let X be a topological space. Let \(\mathcal{C}_{b}(\mathrm{X})\) be the commutative stellar algebra of bounded continuous complex-valued functions on X.
\end{enumerate}

\begin{enumerate}
\def\labelenumi{\alph{enumi})}
\item
  Let \(\widehat{\mathrm{X}} = \mathrm{X}(\mathcal{C}_b(\mathrm{X}))\) ; the space \(\widehat{\mathrm{X}}\) is a compact space such that \(\mathcal{C}(\widehat{\mathrm{X}})\) is identified with \(\mathcal{C}_b(\mathrm{X})\) .
\item
  The map \(j: X \to \widehat{X}\) which sends x to the character \(f \mapsto f(x)\) is continuous.
\item
  If X is completely regular (TG, IX, p.~8, def. 4), the map j is a homeomorphism of X onto a dense subspace of \(\widehat{X}\) (To show that the image of j is dense in \(\widehat{X}\) , show that a continuous complex function on \(\widehat{X}\) vanishing on \(j(\mathrm{X})\) is identically zero.)
\item
  The space \(\widehat{X}\) is identified with the Stone–Čech compactification of X (TG, IX, p.~9).
\end{enumerate}

\begin{enumerate}
\def\labelenumi{\arabic{enumi})}
\setcounter{enumi}{16}
\item
  Let A be a C*-algebra and \(x \in \mathrm{A}\) . If \(x = a - b\) with \(a, b\) positive and \(ab = ba = 0\) , then \(a = x^{+}\) and \(b = x^{-}\) .
\item
  Let A be a separable stellar algebra. There exists a sequence \((e_{n})_{n\in\mathbf{N}}\) into \(A_{+}\) such that the associated elementary filter (TG, I, p.~42, def. 7) is an increasing approximate unit of A.
\item
  Let A be a stella algebra. An element \(x\) is said to be systematically positive if \(x \in \mathrm{A}_{+}\) and if \(x\mathrm{A}x\) is dense in A.
\end{enumerate}

\begin{enumerate}
\def\labelenumi{\alph{enumi})}
\item
  Show that a commutative stellar algebra A contains a systematically positive element if, and only if, there exists a locally compact topological space X, countable at infinity, such that A is isomorphic to \(\mathcal{C}_{0}(X)\) .
\item
  If A is unital, then an element x of A is strictly positive if and only if x is positive and invertible.
\item
  Let \(x \in A\) an element that is systematically positive such that \(\|x\| \leqslant 1\) Prove that the sequence \((x^{1/n})_{n \geqslant 1}\) is an increasing approximate unit in A.
\item
  A stellate algebra A contains a systematically positive element if and only if there exists a sequence \((x_{n})_{n\geqslant1}\) in A, which forms an approximate identity for A.
\end{enumerate}

\begin{enumerate}
\def\labelenumi{\arabic{enumi})}
\setcounter{enumi}{19}
\item
  Let A be a C*-algebra, I a closed two-sided ideal of A, and B a closed involutive subalgebra of A. Then B+I is a closed involutive subalgebra of A, and there exists an isometric isomorphism \(B + I/I \rightarrow B/B \cap I\) .
\item
  Let S be the closed unit disk in C. Let A be the involutive Banach algebra of continuous functions from S to C that are holomorphic in the open unit disk, where the involution is given by \(f^{*}(z) = \overline{f(\overline{z})}\) There exists a Hermitian element of A whose spectrum is not contained in R.
\item
  Let \(\mathrm{A} = \mathcal{C}([0,1])\) Show that the ideal I of A generated by the identity function of {[}0, 1{]} is not closed.
\item
  Let A be a unital C*-algebra. We denote by \(A_{+}^{\prime}\) the vector space of linear forms f on A, not necessarily continuous, such that
\end{enumerate}

\[
f (\mathrm{A} _ {+}) \subset [ 0, + \infty [
\]

(“positive linear forms”).

\begin{enumerate}
\def\labelenumi{\alph{enumi})}
\item
  Let \(f \in A_{+}^{\prime}\) . For every sequence \((x_{n})_{n \in \mathbf{N}}\) into \(A_{+}^{\leqslant 1}\) , the family \((f(x_{n}))\) is bounded. (Show that the series \(\sum_{i} f(x_{i})y_{i}\) converges for every positive family \((y_{i})\) into \(\ell^{1}(\mathbf{N})\) .)
\item
  The linear map f is continuous.
\item
  A linear form \(f\) belongs to \(\mathrm{A}_{+}^{\prime}\) if and only if \(f\) is continuous and \(\| f \| = f(1)\) .
\item
  Let \(x \in A_{h}\) a Hermitian element. There exists a positive linear form \(f \in A_{+}^{\prime}\) such that \(|f(x)| = \|x\|\) (Use Exercise 7 of I, p.~174).
\end{enumerate}

\begin{enumerate}
\def\labelenumi{\arabic{enumi})}
\setcounter{enumi}{23}
\tightlist
\item
  Let A be a unital Banach algebra over C. Suppose that \(x \mapsto x^{*}\) and \(x \mapsto x^{\circ}\) are involutions on A such that A, equipped with each one, is a stellar algebra. We denote by \(A_{1}\) and \(A_{2}\) these two stellar algebras.
\end{enumerate}

\begin{enumerate}
\def\labelenumi{\alph{enumi})}
\item
  Using the notation from Exercise 23, show that \(\mathrm{A}_{1,+}^{\prime} = \mathrm{A}_{2,+}^{\prime}\) . One denotes by \(\mathrm{A}_{+}^{\prime}\) this set.
\item
  Let \(x \in A_{1,h}\) an element such that \(x^{*} = x\) . Show that for every \(f \in A_{+}^{\prime}\) , one has \(f(i(x - x^{\circ})) = 0\) . Deduce that \(x \in A_{2,h}\) (Use part (b) of Exercise 23).
\item
  Conclude that \(A_{1} = A_{2}\) , and that a complex Banach algebra admits at most one involution making it a C*-algebra.
\end{enumerate}

\begin{enumerate}
\def\labelenumi{\arabic{enumi})}
\setcounter{enumi}{24}
\tightlist
\item
  Let G be a group. We denote by C{[}G{]} the group algebra of G over C. It is a subspace of the Hilbert space \(\ell^{2}(\mathrm{G})\) . For g in G, we denote by {[}g{]} the element of C{[}G{]} corresponding to g. The elements {[}g{]}, for \(g \in G\) , form an orthonormal basis of \(\ell^{2}(G)\) and C{[}G{]} is dense in \(\ell^{2}(G)\) .
\end{enumerate}

\begin{enumerate}
\def\labelenumi{\alph{enumi})}
\tightlist
\item
  The map
\end{enumerate}

\[
\sum_ {g} \alpha_ {g} [ g ] \mapsto \sum_ {g} \overline {{\alpha}} _ {g} [ g ^ {- 1} ]
\]

is an involution on \(\mathbf{C}[\mathrm{G}]\) .

\begin{enumerate}
\def\labelenumi{\alph{enumi})}
\setcounter{enumi}{1}
\tightlist
\item
  For \(g \in G\) , let \(\pi(g)\) the map \(x \mapsto gx\) of C{[}G{]}. The map which associates \(\pi(g)\) to g extends to a continuous and injective morphism of involutive algebras from C{[}G{]} into \(\mathcal{L}(\ell^{2}(G))\) .
\end{enumerate}

We denote \(\operatorname{Stell}_{r}(\mathbf{G})\) , and the reduced C*-algebra of G is called the closure in \(\mathcal{L}(\ell^{2}(\mathbf{G}))\) of the image of the map \(\pi\) . It is a C*-algebra. We identify \(\mathbf{C}[\mathbf{G}]\) with its image under \(\pi\) into \(\operatorname{Stell}_{r}(\mathbf{G})\) .

\begin{enumerate}
\def\labelenumi{\alph{enumi})}
\setcounter{enumi}{2}
\tightlist
\item
  The linear trace map \(\operatorname{Tr} \colon \mathbf{C}[\mathrm{G}] \to \mathbf{C}\) such that \(\operatorname{Tr}(1) = 1\) and \(\operatorname{Tr}(g) = 0\) for \(g \neq 1\) extends to a continuous linear map \(\operatorname{Stell}_r(\mathrm{G}) \to \mathbf{C}\) .
\end{enumerate}

\begin{enumerate}
\def\labelenumi{\alph{enumi})}
\setcounter{enumi}{3}
\tightlist
\item
  We have \(\operatorname{Tr}(xy) = \operatorname{Tr}(yx)\) for all \(x\) and \(y\) into \(\operatorname{Stell}_r(G) \) .
\end{enumerate}

\begin{enumerate}
\def\labelenumi{\alph{enumi})}
\setcounter{enumi}{4}
\item
  We have \(\operatorname{Tr}(xx^{*}) \geqslant 0\) for all \(x\) into \(\operatorname{Stell}_r(G)\) , and \(\operatorname{Tr}(xx^{*}) = 0\) if and only if \(x = 0\) .
\item
  Let e be an idempotent element of C{[}G{]} such that \(e \neq 0\) and \(e \neq 1\) . Show that \(0 < \operatorname{Tr}(e) < 1\) . (Use Exercise 3).
\item
  Let e be an idempotent element of Z{[}G{]}. Show that e = 0 or e = 1. ("Kaplansky's theorem").
\item
  Let x be an element of C{[}G{]}. Then x is invertible if and only if it is left-invertible, if and only if it is right-invertible.
\item
  Let n be a strictly positive integer. A matrix x of type \((n, n)\) with coefficients in C{[}G{]} is invertible if and only if it is left-invertible, if and only if it is right-invertible.
\end{enumerate}

\begin{enumerate}
\def\labelenumi{\arabic{enumi})}
\setcounter{enumi}{25}
\item
  Let A be a unital C*-algebra. Let \(a\) and \(b\) be normal elements of A and \(x \in \mathrm{A}^*\) such that \(b = xax^{-1}\) . There exists a unitary element \(u\) of A such that \(b = uau^*\) . (Use the preceding exercise and the polar decomposition of \(x\) , noting that \(|x|\) belongs to the bicommutant of \(x\) ).
\item
  Let A be a unital C*-algebra, and let a, b, and c be elements of A. Suppose that b and c are normal. If ac = cb, then we have \(f(a)c = cf(b)\) for every continuous function f on the union of the spectrum of a and the spectrum of b.
\item
  Let A be a C*-algebra. Let \(a\) and \(b\) be positive elements of A such that \(ab \neq 0\) . Then \(\mathrm{Sp}_{\mathrm{A}}(ab)\) is not reduced to 0.
\item
  \begin{enumerate}
  \def\labelenumii{\alph{enumii})}
  \tightlist
  \item
    Let G be a locally compact topological group. The Banach algebra L \(^{1}\) (G) has an approximate identity. (Use INT, VIII, § 2, no. 7, lemma 4).
  \end{enumerate}
\end{enumerate}

\begin{enumerate}
\def\labelenumi{\alph{enumi})}
\setcounter{enumi}{1}
\item
  Let A be a Banach algebra admitting an approximate identity. For every character \(\chi\in\mathsf{X}(\mathrm{A})\) , one has \(\|\chi\|=1\) .
\item
  The Banach subalgebra of L \(^{1}\) (Z) formed by the sequences ( \(x_{n}\) such that \(x_{n} = 0\) if \(n < 0\) has no approximate identity. (Construct a nonzero character of norm \(< 1\) .)
\end{enumerate}

\begin{enumerate}
\def\labelenumi{\alph{enumi})}
\setcounter{enumi}{3}
\tightlist
\item
  Let G be the free group generated by two elements \(a\) and \(b\) and \(\mu\) a Haar measure on G. Let I be the closed ideal of \(\mathrm{L}^1 (\mathrm{G})\) consisting of the \(f\in \mathrm{L}^{1}(\mathrm{G})\) such that \(\mu (f) = 0\) . Let \(\xi = e^{2i\pi /3}\) . Let \(x = \varepsilon_{e} + \xi \varepsilon_{a} + \xi^{-1}\varepsilon_{b}\in \mathrm{L}^{1}(\mathrm{G})\) . For every \(f\in \mathrm{L}^{1}(\mathrm{G})\) , one has \(\| x*f - x\| \geqslant 1\) The Banach algebra I has no approximate identity.
\end{enumerate}

\begin{enumerate}
\def\labelenumi{¶ \arabic{enumi})}
\setcounter{enumi}{29}
\tightlist
\item
  Let A be a Banach algebra with an approximate identity. Denote by \(\widetilde{\mathsf{A}}\) the Banach algebra obtained from A by adjoining an identity element, denoted by 1.
\end{enumerate}

\begin{enumerate}
\def\labelenumi{\alph{enumi})}
\item
  Let \(x \in \mathrm{A}\) . Then \(x\) belongs to the closed left ideal generated by \(x\) .
\item
  For every \(\varepsilon > 0\) and every finite family \((x_i)_{i \in \mathrm{I}}\) of elements of \(\mathrm{A}\) , there exists \(e \in \mathrm{A}\) such that \(\| e \| \leqslant 1\) and \(\| ex_i - x_i \| < \varepsilon\) for all \(i \in \mathrm{I}\) .
\item
  Let \(x \in A\) , \(\varepsilon > 0\) and \(0 < \gamma \leqslant 1/8\) . There exists \(\eta > 0\) such that for all \(e \in A\) satisfying \(\|e\| \leqslant 1\) and \(\|ex - x\| < \eta\) , then \(\alpha = (1 - \gamma) \cdot 1 + \gamma e\) is invertible in \(\widetilde{A}\) and \(\|x - \alpha^{-1}x\| < \varepsilon\) .
\end{enumerate}

\begin{enumerate}
\def\labelenumi{\alph{enumi})}
\setcounter{enumi}{3}
\tightlist
\item
  Let \(x \in \mathrm{A}\) , \(\varepsilon > 0\) , \(\delta > 0\) and \(0 < \gamma \leqslant 1/8\) . There exists a sequence \((e_n)_{n \geqslant 1}\) in A having the following properties:

\begin{enumerate}
\def\labelenumi{(\roman{enumi})}
\item
  We have \(\|e_{n}\| < 2\) for all \(n \geqslant 1\) ;
\item
  The element
\end{enumerate}

\[
y _ {n} = \sum_ {k = 1} ^ {n} \gamma (1 - \gamma) ^ {k - 1} e _ {k} + (1 - \gamma) ^ {n} \cdot 1
\]

is invertible in \(\widetilde{\mathrm{A}}\) for all \(n\geqslant 1\) and

\[
\left\| y _ {n} ^ {- 1} x - y _ {n - 1} ^ {- 1} x \right\| <   \delta 2 ^ {- n} \quad (n \geqslant 2), \quad \left\| y _ {1} ^ {- 1} x - x \right\| <   \delta / 2.
\]

(Assuming that \((e_{1},\ldots,e_{n})\) have been defined, consider an element \(e_{n+1}\) arbitrary of norm \textless{} 2; observe that

\[
y _ {n + 1} = ((1 - \gamma) \cdot 1 + e _ {n + 1}) \left(\sum_ {k = 1} ^ {n} \gamma (1 - \gamma) ^ {k - 1} e _ {k} ^ {\prime} + (1 - \gamma) ^ {n} \cdot 1\right),
\]

where \(e_k' = ((1 - \gamma) \cdot 1 + e_{n+1})^{-1} e_k\) , and determine \(e_{n+1}\) satisfying the required conditions by exploiting the fact that the set of invertible elements of \(\widetilde{\mathsf{A}}\) is open, and that the map \(a \mapsto a^{-1}\) is continuous.)
\end{enumerate}

\begin{enumerate}
\def\labelenumi{\alph{enumi})}
\setcounter{enumi}{4}
\tightlist
\item
  For every \(x \in A\) , there exist y and z in A such that x = yz ("Cohen's factorization theorem"). Let \(\delta > 0\) . One may assume that z belongs to the closed left ideal generated by \(x\) and that \(\|y-z\|<\delta\) . (Consider a sequence \((e_{n})_{n\geqslant1}\) as above and define \(y_{n}\) as before; the sequence \(z_{n}=y_{n}^{-1}x\) converges to an element \(z\in\mathrm{A}\) , and \(z_{n}\) belongs to the closed left ideal generated by \(x\) ; the sequence \((y_{n})\) converges to \(y\in\mathrm{A}\) and \(x=yz\) .)
\end{enumerate}

\begin{enumerate}
\def\labelenumi{\arabic{enumi})}
\setcounter{enumi}{30}
\tightlist
\item
  Let A be the Banach space of complex integrable functions on {[}0, 1{]} with respect to Lebesgue measure. Set
\end{enumerate}

\[
(f \cdot g) (t) = \int_ {0} ^ {t} f (t - x) g (x) d x
\]

for \(t \in [0, 1]\) .

\begin{enumerate}
\def\labelenumi{\alph{enumi})}
\item
  The function \(f \cdot g\) is defined almost everywhere and is an element of A.
\item
  Equipped with the operation \((f,g)\mapsto f\cdot g\) , the space A is a commutative Banach algebra.
\item
  Let \(x \in \mathrm{A}\) be the constant function equal to 1. The function \(x^n\) is
\end{enumerate}

\[
t \mapsto \frac {t ^ {n - 1}}{(n - 1) !}.
\]

\begin{enumerate}
\def\labelenumi{\alph{enumi})}
\setcounter{enumi}{3}
\item
  The algebra A is topologically generated by x, the element x is quasi-nilpotent, and A is equal to its radical.
\item
  The algebra A admits an approximate identity.
\item
  The algebra A has no maximal ideal, regular or not, closed or not (use the preceding exercise and exercise 1 of I, p.~166). In particular, X(A) is empty, hence compact.
\end{enumerate}

\begin{enumerate}
\def\labelenumi{\arabic{enumi})}
\setcounter{enumi}{31}
\tightlist
\item
  Let G be a group. One says that G has the Powers property if, for every \(n \in N\) and every finite subset F of \(G = \{e\}\) , there exists a partition of G into two subsets A and B and a family \((g_i)_{1 \leqslant i \leqslant n}\) of elements of G such that
\end{enumerate}

\begin{enumerate}
\def\labelenumi{(\roman{enumi})}
\item
  for all \(g \in \mathrm{F}\) , the sets A and gA are disjoint;
\item
  for all distinct integers \(j\) and \(k\) between 1 and \(n\) , the sets \(g_{j}\mathrm{B}\) and \(g_{k}\mathrm{B}\) are disjoint.
\end{enumerate}

Let G be a group possessing the Powers property. Let I be a nonzero closed two-sided ideal of the reduced C*-algebra \(\text{Stell}_{r}(\text{G})\) (exercise 25). Let \(\pi: C[G] \to \text{Stell}_{r}(G)\) be the canonical representation.

\begin{enumerate}
\def\labelenumi{\alph{enumi})}
\tightlist
\item
  There exists a family \((\lambda_g)_{g\in \mathrm{G - }\{e\}}\) into \(\mathbf{C}\) such that
\end{enumerate}

\[
x = \sum_ {g \in \mathrm{G} - \{e \}} \lambda_ {g} \pi (g)
\]

is a hermitian element of \(\operatorname{Stell}_{r}(\mathrm{G})\) such that \(y = 1 + x\) belongs to I.

\begin{enumerate}
\def\labelenumi{\alph{enumi})}
\setcounter{enumi}{1}
\tightlist
\item
  Let \(\varepsilon > 0\) such that \(\varepsilon < 1\) . There exists a finite subset F of G = \{e\} and an integer \(n \geqslant 1\) such that the element
\end{enumerate}

\[
x ^ {\prime} = \sum_ {g \in \mathrm{F}} \lambda_ {g} \pi (g)
\]

satisfies \(\|x' - x\| < \varepsilon\) and \(\|x'\| < \frac{\sqrt{n}}{2}(1 - \varepsilon)\) .

\begin{enumerate}
\def\labelenumi{\alph{enumi})}
\setcounter{enumi}{2}
\tightlist
\item
  Let \(G = A \cup B\) and \((g_{i})_{1 \leqslant i \leqslant n}\) satisfying the conditions of the Powers property for the finite subset F and the integer n.~Set
\end{enumerate}

\[
y ^ {\prime} = \frac {1}{n} \sum_ {i = 1} ^ {n} \pi (g _ {i}) x ^ {\prime} \pi (g _ {i} ^ {- 1}).
\]

We have \(\|y'\| < 1 - \varepsilon\) . (Let \(p_i\) be the orthogonal projection of \(\ell^{2}(\mathrm{G})\) on \(\ell^{2}(g_{i}\mathrm{B})\) ; verify that

\[
(1 - p _ {i}) \pi (g _ {i}) x ^ {\prime} \pi (g _ {i} ^ {- 1}) (1 - p _ {i}) = 0
\]

for \(1 \leqslant i \leqslant n\) , and deduce that

\[
1 + y ^ {\prime} = 1 + \frac {1}{n} \sum_ {i = 1} ^ {n} p _ {i} \pi (g _ {i}) x ^ {\prime} \pi (g _ {i} ^ {- 1}) + \left(\frac {1}{n} \sum_ {i = 1} ^ {n} p _ {i} \pi (g _ {i}) x ^ {\prime} \pi (g _ {i} ^ {- 1}) (1 - p _ {i})\right) ^ {*},
\]

then note that the endomorphisms \(p_{i}\pi(g_{i})x^{\prime}\pi(g_{i}^{-1})\) have pairwise orthogonal images to estimate their norms.)

\begin{enumerate}
\def\labelenumi{\alph{enumi})}
\setcounter{enumi}{3}
\tightlist
\item
  The ideal I contains the element
\end{enumerate}

\[
1 + \frac {1}{n} \sum_ {i = 1} ^ {n} \pi (g _ {i}) x \pi (g _ {i} ^ {- 1}),
\]

and this element is invertible. Consequently, the stellar algebra \(\text{Stell}_{r}(\text{G})\) contains no closed two-sided ideal other than \(\{0\}\) and of \(\text{Stell}_{r}(\text{G})\) .

\begin{enumerate}
\def\labelenumi{\alph{enumi})}
\setcounter{enumi}{4}
\tightlist
\item
  Let G be a noncommutative free group. Then G has the Powers property.
\end{enumerate}
% semantic-fragments: legacy-input-seam
\relax{}
